\section{Related Work}

AI-based simulation has emerged as an effective method for developing interpersonal soft skills across domains. \citet{yang2024} provide a unifying framework for these systems, distinguishing an \textit{AI Partner} component---which simulates a human interlocutor for deliberate practice---from an \textit{AI Mentor} component, which delivers structured feedback based on domain expertise. \citet{shaikh2024} developed Rehearsal, a conflict resolution training system grounded in the Interests-Rights-Power (IRP) framework. Its AI Partner simulates a conflict counterpart across eight IRP-grounded strategies; its AI Mentor delivers in-context, turn-by-turn feedback---presenting counterfactual alternative messages and prompting users to identify the strategy employed. A controlled experiment ($N=40$) found that simulated practice significantly increased participants' use of cooperative strategies and reduced competitive strategies in a live conflict with a blinded confederate ($p < 0.05$), while producing no significant improvement on knowledge recall, suggesting AI simulation builds applied skill rather than declarative knowledge.

\citet{daryanto2025} developed Conversate, combining an interview roleplay simulation with a post-session feedback agent. After the simulation, the system identifies weak areas in the transcript; users select specific moments, write a self-assessment note, and then initiate a chat-based feedback session with the AI focused on those moments, grounded in dialogic feedback theory---which structures feedback as a two-way conversation where learners actively seek meaning from it. The authors observe that LLM sycophancy---the model's tendency to concede when users challenge its assessment---emerges as a recurring limitation; interview skill improvement was not measured.

Negotiation, as one such interpersonal skill domain, has similarly emerged as a promising target for AI-based simulation training \citep{dannenmann2022, rottner2024}. \citet{dinnar2021} note that negotiation students learn best when they reflect on their own performance, receive personalized feedback, and integrate both into their learning. \citet{dannenmann2022} developed a computer-based simulation in which trainees negotiated with NLP-powered virtual agents before receiving a post-session overview of their negotiation style and outcomes from an AI coach; a between-subjects experiment ($N=27$) found significantly greater gains in evaluator-assessed negotiation performance than traditional role-play, suggesting that technology-supported approaches can outperform conventional methods. \citet{rottner2024} developed AdVentures, a two-round negotiation exercise in which students practice with an AI counterpart, receive personalized AI feedback between rounds, and then renegotiate; an initial pilot study ($N=53$) found a 10\% improvement in negotiation outcomes across rounds.

\citet{shea2024} developed ACE, a system in which users practice with an LLM-based negotiation agent and then receive a written coaching report identifying tactical errors, providing targeted feedback, and suggesting utterance revisions; a controlled experiment ($N=374$) found that only ACE's annotation-based feedback improved negotiation outcomes, while generic AI-generated feedback performed no better than no feedback. More recently, \citet{lirong2025} evaluated a generative AI negotiation coach featuring dynamically adaptive opponents capable of assuming various negotiation styles and providing real-time personalized feedback, finding significantly greater gains in strategic preparation, strategy, and value creation than traditional peer-to-peer role-play, though the control group showed significantly stronger development in communication skills.

Yet none of these systems closes \citeauthor{dinnar2021}'s \citeyearpar{dinnar2021} learning cycle---the self-reflection step remains largely unaddressed. As \citet{rudolph2007} argue, identifying errors is insufficient: lasting improvement requires bringing the underlying cognitive frames that drive behavior to the surface. Simulation debriefing is a structured conversational intervention designed to surface those frames---distinct from the written feedback reports that dominate prior systems.

\citet{fanning2007} identify debriefing as the most critical component of the simulation experience---``crucial to the learning process.'' Emerging work has begun applying AI to facilitate debriefing. \citet{hong2025} used GPT-4o to generate debrief scripts to guide human facilitators during pediatric simulations, with both facilitators and learners reporting strong enthusiasm and perceived improvements in debriefing organization. \citet{gonzalez2025} deployed an AI system that walked nursing students through a fixed sequence of reflection prompts and analyzed their spoken responses, finding no significant correlation between time in debrief and simulation performance; the authors conclude that AI-facilitated debriefing ``may not replace the depth of human-facilitated sessions.''

Evangelou et al.~\citep{evangelou2025, evangelou2026} applied AI debriefing to VR-based counselor training. A GenAI chatbot led students through a structured three-phase protocol---opening with emotional reactions to the simulation, moving to analysis of specific techniques practiced, and closing with transfer to future practice---progressing through each phase in a fixed sequence. The 2025 qualitative study identified sycophantic tendencies as a recurring limitation; the 2026 controlled study found no significant differences in cognitive load between chatbot- and human-led debriefings. Notably, Debriefing with Good Judgment \citep{rudolph2007} offers a theoretically grounded conversational method, combining explicit advocacy with genuine inquiry, that none of these systems has applied.

No prior negotiation training system combines AI-based practice with a post-session dialogic AI coach implementing a theoretically grounded reflective methodology. This project addresses both gaps by combining an AI negotiation simulator with an autonomous post-session debrief agent that diagnoses negotiation behavior against \textit{Getting to Yes} principles \citep{fisher1981} and guides reflection through Debriefing with Good Judgment, evaluated through a controlled experiment.
